\section{Discussion}
\label{sec:discussion}

\subsection{Key Findings}

\textbf{Verbosity is the primary annotation drift channel.} Across all three experiments, verbosity ($\beta_L$) is the only stylistic feature that consistently achieves statistical significance. Verbosity is the most perceptually salient feature of AI-generated text, and automation bias theory predicts that annotators under cognitive load will rely on the most accessible surface signal. The sign reversal --- from penalizing to rewarding verbosity --- is consistent with annotators progressively adopting a ``more detailed $=$ better quality'' heuristic that tracks AI output style.

The subthreshold results for $\beta_H$ and $\beta_S$ are best interpreted as reflecting limits of the current proxy design. All three deltas are positive across 160,800 annotation decisions --- the directional consistency is informative.

\textbf{The AI-typicality projection confirms directional specificity.} The geometric projection result ($\beta_\text{exposure} = 0.041$, $p = 2.05\times10^{-5}$; placebo $p = 0.48$) establishes that the preference drift is not in an arbitrary stylistic direction --- it is specifically aligned with the axis from human-typical to AI-typical text in sentence embedding space.

\textbf{Null results specify the data frontier.} The inability to test ambiguity moderation or within-annotator dose-response is a precise characterization of what existing RLHF datasets can and cannot support.

\subsection{Limitations}

\textbf{L1: No genuine temporal metadata in HH-RLHF.} Round stratification uses index-based partitioning. The observed $\beta_L$ reversal could reflect between-cohort composition differences rather than within-annotator adaptation. \emph{Framing:} ``We present our findings as population-level directional evidence pending datasets with genuine annotation timestamps.''

\textbf{L2: WebGPT within-annotator design collapsed to between-group.} Worker IDs absent from public release; tercile proxy cannot rule out selection effects. Discriminant validity is confirmed; between-group direction is consistent with H-M2.

\textbf{L3: AAI composite only 2/3 components validated.} The behavioral divergence component (H-M3: reward model comparison) and benchmark degradation component (H-M4) are not executed. H-M3 is fully specified and uses validated code components from H-M2.

\textbf{L4: Topic distribution imbalance ($p = 4\times10^{-275}$).} May confound $\beta_H$ and $\beta_S$ estimates but is less likely to affect $\beta_L$ (response length is less topic-specific).

\subsection{Broader Impact}

The AAI framework provides a practical, low-cost quality gate: if round-stratified coefficient comparison reveals $\beta_L$ sign change across annotation phases, practitioners should investigate drift contamination before reward model training. This requires only preference labels and response text --- no additional annotation overhead.

Our findings suggest that multi-round RLHF annotation datasets may encode temporal non-stationarity that is invisible to standard pairwise accuracy metrics. Reward models trained on later-round labels may have subtly different optimization targets than those trained on earlier-round labels. We do not claim that annotation drift is present in all RLHF datasets or necessarily causes harmful outcomes; our contribution is the measurement instrument that makes this question empirically tractable.
